\documentclass[11pt]{article}

\usepackage[margin=1in]{geometry}
\usepackage[T1]{fontenc}
\usepackage[utf8]{inputenc}
\usepackage{lmodern}
\usepackage{setspace}
\usepackage{hyperref}
\usepackage{booktabs}
\usepackage{csquotes}
\usepackage[numbers,sort&compress]{natbib}

\title{From Behavioral State Vocabularies to Distributed Embodiment:\\
SOFTSIM and GA144 as an Executable Framework}
\author{Christopher A. Tucker}
\date{}

\begin{document}

\maketitle

\begin{abstract}
This paper is the implementation-oriented third step in a sequence of
related studies. A prior historical paper argued that an older line of
machine-intelligence thought remains architecturally useful when read in
terms of sensing, state, control, action, and adaptation. A second paper
showed that Sony's \texttt{ERS-111} \texttt{R-CODE} sample corpus can be
read through a compact behavioral state vocabulary rather than as a set of
unrelated scripts. The present paper asks what follows if that vocabulary is
treated not only as an analytic description, but as a candidate
implementation grammar. Its main claim is that the GreenArrays toolchain,
especially \texttt{arrayForth}, boot-descriptor loading, \texttt{SOFTSIM},
and physical \texttt{GA144} deployment, provides a continuous path from
behavioral decomposition to executable distributed embodiment. On this
reading, behavior diagrams are neither loose metaphors nor final programs.
They are intermediate control descriptions that can be mapped onto
communicating \texttt{F18A} nodes, initialized reproducibly, exercised
functionally in simulation, and then carried toward hardware realization.
The paper positions this workflow as a concrete framework for studying
fast-reactive and slow-supervisory embodied control without collapsing both
layers into one monolithic controller.
\end{abstract}

\onehalfspacing

\section{Introduction}

This paper is written as the third in a sequence. The first paper argued
that a historically underused architecture of machine intelligence can be
recovered by reading symbolic logic, circuitry, control, and embodied
behavior together rather than apart \citep{tucker2026berkeley}. The second
paper moved from historical reconstruction to preserved robotic software
and argued that Sony's \texttt{ERS-111} \texttt{R-CODE} sample set is
organized by a compact behavioral state vocabulary centered on
initialization, sensing, action, synchronization, and recovery
\citep{tucker2026behavioral}. The present paper takes the next step. It
asks how such a vocabulary can be made executable on a distributed
substrate instead of remaining a descriptive layer over legacy code.

The relation among the three papers is therefore architectural rather than
merely chronological. The first paper supplies the engineering hypothesis:
embodied machine intelligence can be organized through explicit relations
among sensing, state or memory, control, action, recovery, and adaptation.
The second paper supplies the intermediate behavioral representation: a
reusable vocabulary of states and transitions extracted from a preserved
robot corpus. The present paper treats that representation as something to
be instantiated. Its task is to ask whether the vocabulary can be mapped
onto executable node-local processes, initialized reproducibly, inspected in
simulation, and carried toward a distributed hardware substrate.

That question matters because much current robotics work still separates
high-level cognition from low-level embodied viability too sharply. Systems
may be rich in perception, planning, or language conditioning while
remaining comparatively weak in persistent local organization: postural
maintenance, contact response, local recovery, and bounded reactive loops
that keep a physical agent workable in the world while larger tasks are
being pursued. One common way of talking about that split is through
fast/slow or \emph{System 1} / \emph{System 2} language, where rapid
background response and slower supervisory interpretation are treated as
distinct but cooperating cadences \citep{kahneman2011thinking,wang2025hume}.

The aim of this paper is not to claim that GreenArrays hardware solves
contemporary robotics in general. It is narrower and more practical. The
claim is that the \texttt{F18A} / \texttt{arrayForth} /
\texttt{SOFTSIM} / \texttt{GA144} stack offers a usable experimental path
for implementing distributed embodied control in explicit small processes.
That makes it a good setting in which to test whether behavioral
decompositions derived from robot code can be turned into node-local loops,
inter-node handoffs, and slower supervisory interventions without being
absorbed into a single centralized controller.

\section{Research Question and Thesis}

The guiding question is how a behavioral state vocabulary should function
after analysis. Is it only a retrospective interpretive device for reading
robot scripts, or can it serve as an intermediate implementation layer that
connects behavior design to executable distributed systems?

The thesis of this paper is that the GreenArrays toolchain supports the
stronger reading. In this environment, a behavior description can be mapped
onto concrete communicating \texttt{F18A} nodes, assembled into persistent
per-node images, initialized declaratively through boot-descriptor
mechanisms, executed functionally in \texttt{SOFTSIM}, and then carried to
physical \texttt{GA144} hardware \citep{greenarrays2011db001,greenarrays2011db013}. The
result is not yet a finished robotic architecture. It is an executable
framework for studying how fast reactive loops, local state, and slower
supervisory control might be separated, coordinated, inspected, and
eventually embodied in hardware. In that sense, the paper does not treat
the earlier historical and corpus studies as proof of the implementation
claim. It treats them as the conceptual and representational prerequisites
for a concrete test of distributed embodiment. The proof burden remains
execution: node allocation, load structure, cadence separation, behavioral
handoff, simulation evidence, and eventual hardware evidence must be shown
rather than inferred from the earlier analyses.

\section{From Behavioral Vocabulary to Implementation Grammar}

The earlier \texttt{R-CODE} study argued that apparently different Sony
sample routines reuse a small recurrent vocabulary of control forms:
startup regularization, sensing, conditional branching, iterative action,
synchronization, and recovery \citep{tucker2026behavioral}. That finding
has two possible interpretations. The weaker one is taxonomic: the
vocabulary merely groups legacy scripts into comparable families. The
stronger one is constructive: the vocabulary identifies a set of behavior
blocks that can be re-instantiated on another substrate.

This paper adopts the stronger interpretation. In that reading, a behavior
diagram is not yet machine code, but it is more than a sketch. It specifies
roles that an implementation must realize: persistent loops, transition
conditions, local memory, recovery paths, and coordination points. Such a
description is therefore well suited to a substrate whose basic unit is not
a single giant process but a collection of communicating node-local
programs.

That is where the GreenArrays environment becomes relevant. The \texttt{F18A}
architecture is already distributed in form. It consists of many small
communicating nodes with explicit port semantics, node-local memory, and
fine-grained execution \citep{greenarrays2011db001}. The programming model
described by \texttt{arrayForth} then reinforces that decomposition by
treating code as persistent per-node object images rather than as one flat
binary \citep{greenarrays2011db013}. What the earlier paper recovered at
the level of behavioral analysis thus meets a substrate that is already
organized for decomposition.

\section{The GreenArrays Execution Path}

The practical value of the GreenArrays stack lies in the continuity of its
workflow. The manuals and local notes together support a five-stage path:
behavioral model, node allocation, per-node assembly, declarative
initialization, and executable realization.

The first stage is the behavior description itself. Here the relevant object
is a compact control representation: a loop, a sensing condition, a
supervisory trigger, a recovery branch, or a state-maintaining routine. The
second stage assigns these roles to concrete nodes and communication paths.
This is the point at which a state machine or behavior ladder becomes a
distributed organization rather than a conceptual diagram.

The third stage is per-node implementation. \texttt{arrayForth} Chapter 4
describes an incremental environment in which each node has a persistent
object bin containing RAM image, ROM image, and labels
\citep{greenarrays2011db013}. That matters because it makes implementation
structurally congruent with the distributed behavior idea. One does not
begin from a monolithic application image and then discover decomposition
afterward. The decomposition is present in the object model from the start.

The fourth stage is declarative initialization. Chapter 6's boot-descriptor
machinery provides a way to specify how many nodes are touched, how memory
and registers are initialized, and how a node ensemble is brought into
execution-ready form \citep{greenarrays2011db013}. This is not just a
loading convenience. It is part of the argument of the paper. If a behavior
architecture is to be studied seriously, it must be instantiated
reproducibly. The boot-descriptor layer supplies that reproducible
instantiation path.

The fifth stage is executable realization. \texttt{SOFTSIM} provides a
functional environment in which the resulting node program can be loaded,
stepped, and inspected, while \texttt{GA144} provides the physical
distributed substrate on which the same style of program can be manifested.
The important constraint is that the simulator should be treated as a
functional witness, not as a timing-faithful proxy for the hardware. It can
validate organization and value-level behavior, while real-chip testing
remains necessary for timing, contention, and physical I/O questions.

\section{Fast and Slow Cadence}

One reason this implementation path is worth isolating is that it offers a
concrete place to study fast and slow control decomposition. Public robotics
literature now contains multiple explicit versions of this pattern, in which
a faster reactive layer cooperates with a slower interpretive or
value-guided layer \citep{wang2025hume}. The GreenArrays stack is relevant
not because it replicates those systems wholesale, but because it offers a
different implementation emphasis: distributed low-latency node processes
instead of one heavy central controller.

In the present framing, fast control corresponds to local reactive loops:
posture maintenance, contact response, bounded actuator logic, sensing, and
recovery. Slow control corresponds to supervisory choice, mode transition,
task redirection, or larger behavioral revision. This is conceptually close
to familiar behavior-based and layered-control traditions
\citep{arkin1998behavior,brooks1991intelligence}, but the present paper adds
an implementation claim. It argues that the split can be explored on a
distributed asynchronous substrate whose programming and loading machinery
already respects node locality.

This has methodological value. It lets a researcher ask not only whether a
behavior architecture is intelligible on paper, but whether its parts can be
given separate executable homes, separate initialization conditions, and
separate inspection points. That makes the architecture more testable. One
can begin with a minimal one-node reactive loop, then expand to a two-node
handoff, then move toward richer supervisory decomposition while keeping the
relationship among behavior, implementation, and loading visible.

\section{A Minimal Research Program}

The current local notes suggest a practical research program that starts
small and stays traceable. The first step is not a large multi-node robotic
brain. It is a minimal structural mutation of a known-good \texttt{SOFTSIM}
load pattern. That means preserving a confirmed small example, introducing a
controlled one-node variant, verifying that it still loads and steps
correctly, and only then extending the pattern toward a two-node cadence
experiment.

This sequence is methodologically important because it separates three kinds
of error that would otherwise be conflated: syntax or assembly error,
load-and-initialization error, and inter-node coordination error. Once these
are disentangled, the implementation path becomes suitable for more serious
behavior experiments. A reactive loop can be isolated, a supervisory trigger
can be added, and the resulting handoff can be inspected first in
\texttt{SOFTSIM} and later on hardware.

The broader value proposition is architectural rather than benchmark-driven.
\texttt{GA144} is interesting here because it offers a plausible low-energy,
many-node substrate for persistent small embodied loops. The point is not
that current notes prove superior robot performance or timing efficiency.
The point is that the substrate fits the decomposition under study:
background reactive processes that remain active continuously rather than as
occasional subroutines.

\section{Contribution of the Third Paper}

Within the three-paper sequence, the present paper occupies a distinct role.
The first paper recovered an architectural lineage for embodied machine
intelligence. The second showed that a preserved robot-code corpus can be
read through a stable behavioral state vocabulary. The third paper turns
that vocabulary toward execution. Its contribution is to identify a concrete
toolchain in which behavior blocks, node decomposition, initialization, and
functional execution can be studied together rather than in isolation.

That contribution should be stated carefully. The paper does not yet present
a completed distributed robot controller, nor does it claim timing
validation, energy benchmarking, or parity with contemporary foundation-model
robot systems. Its claim is narrower and still useful: behavior diagrams can
serve as an intermediate representation for executable distributed
embodiment, and the GreenArrays stack provides a disciplined environment in
which that claim can be explored incrementally.

\section{Conclusion}

The main result of this paper is a framework claim. A behavioral state
vocabulary need not remain a retrospective analytic lens over preserved robot
code. It can function as an implementation grammar when paired with a
substrate that already supports node-local programming, explicit
communication, reproducible initialization, and functional execution.

Within that framing, the GreenArrays environment supplies a continuous path
from behavior description to distributed realization: node-level
implementation in \texttt{arrayForth}, declarative ensemble setup through
boot descriptors, functional execution in \texttt{SOFTSIM}, and eventual
physical manifestation on \texttt{GA144}. This does not complete the larger
research program, but it makes the next step concrete. Instead of discussing
fast and slow embodied control only at the level of architecture, we can now
study how such a decomposition behaves when given explicit executable form on
a distributed substrate.

\bibliographystyle{unsrtnat}
\bibliography{references}

\end{document}
